\section{Experiments \& Results}
\label{sec:exp}
% [6 pts]  Quantitative results (CLIP score) + qualitative results (output images).
% [6 pts]  Ablation: (1) SDS guidance_scale; (2) SDI inversion_n_steps; (3) SDI sdi_update_interval.
% [2 pts]  All hyperparameters for each task.
%
% Images: copy outputs/{sds,vsd,sdi}/*.png to Figure/{sds,vsd,sdi}/.
% A missing image is drawn as a framed placeholder box showing its expected path.

\input{Figure/figure_page}  % Figs. 2-6 together on one figure page

%-------------------------------------------------------------------------
\subsection{Quantitative Results}
\label{sec:quantitative}

\begin{table}[t]
  \centering
  \footnotesize
  \setlength{\tabcolsep}{4pt}
  \caption{\textbf{CLIP score} ($\uparrow$, ViT-B/32) per prompt.}
  \label{tab:clip}
  % values: outputs/{sds,vsd,sdi}/eval.json
  \begin{tabular}{@{}lccc@{}}
    \toprule
    Prompt & SDS & VSD & SDI \\
    \midrule
    A red bus driving on a desert road & 0.3442 & 0.3528 & \textbf{0.3846} \\
    a boat in a river & 0.2953 & \textbf{0.3034} & 0.2992 \\
    A cabin surrounded by forests & 0.3010 & 0.3228 & \textbf{0.3247} \\
    A church beside a lake & 0.3246 & 0.3224 & \textbf{0.3362} \\
    A villa close to the pool & 0.2935 & \textbf{0.3097} & 0.3006 \\
    A castle next to a river & 0.2959 & \textbf{0.3009} & 0.3008 \\
    A burger on the table & \textbf{0.3047} & 0.2874 & 0.2990 \\
    A dog sitting on grass & \textbf{0.2880} & 0.2688 & 0.2627 \\
    a cat sitting on a table & 0.2989 & 0.2725 & \textbf{0.3021} \\
    A car on the road & \textbf{0.2953} & 0.2699 & 0.2757 \\
    \midrule
    \textbf{Average} & 0.3042 & 0.3011 & \textbf{0.3086} \\
    \bottomrule
  \end{tabular}
\end{table}


\Cref{tab:clip} reports CLIP scores for the ten evaluation prompts.
SDI has the highest average CLIP score (0.3086), followed by SDS (0.3042) and VSD (0.3011).
SDI leads on four prompts, while SDS and VSD each lead on three.

The score ranking only partly agrees with the visual comparison in \cref{fig:qualitative}.
For the burger, SDS scores 0.3047 despite producing a blurrier image than SDI, which scores 0.2990.
The object remains recognizable in SDS, illustrating that stronger image--text similarity does not
necessarily imply sharper details.

%-------------------------------------------------------------------------
\subsection{Qualitative Results}
\label{sec:qualitative}

\Cref{fig:qualitative} shows clear differences in texture and structure among the three methods.
SDS produces recognizable objects and simple compositions, but edges and surfaces are heavily
smoothed.
For example, the burger has a clear silhouette but little texture, and the cabin and castle lose
fine architectural details.

VSD produces sharper local details and more varied colors, but many images contain dense, irregular
textures.
The burger has visible sesame seeds and layered ingredients, while the villa includes vegetation and
pool reflections.
However, mottled patterns cover the burger background, and the villa's walls and roof appear
fragmented.
These artifacts show that additional detail does not always yield coherent surfaces.

SDI generally produces cleaner surfaces and more distinct object boundaries.
The cabin and church have clearer architectural structure, and the burger shows individual
ingredients against a smoother background.

%-------------------------------------------------------------------------
\subsection{Ablation Study}
\label{sec:ablation}

\begin{table}[t]
  \centering
  \footnotesize
  \setlength{\tabcolsep}{4pt}
  \caption{\textbf{Ablation study.} Average CLIP score ($\uparrow$) over the Burger, Cat,
  and Villa prompts. Existing baseline images are reused.}
  \label{tab:ablation}
  \begin{tabular}{@{}lccccc@{}}
    \toprule
    \multicolumn{6}{@{}l}{\textit{(a) SDS: \texttt{guidance\_scale}}} \\
    Value & 7.5    & 15     & 25     & 50     & 100    \\
    CLIP & 0.2426 & 0.2950 & \textbf{0.2991} & 0.2879 & 0.2815 \\
    \midrule
    \multicolumn{6}{@{}l}{\textit{(b) SDI: \texttt{inversion\_n\_steps}}} \\
    Value & 1      & 3      & 5      & 10     & 25     \\
    CLIP & 0.1809 & 0.2942 & 0.2874 & \textbf{0.3006} & 0.2956 \\
    \midrule
    \multicolumn{6}{@{}l}{\textit{(c) SDI: \texttt{sdi\_update\_interval}}} \\
    Value & 25     & 75     & 200    & 500    & 1000   \\
    CLIP & \textbf{0.3006} & 0.2848 & 0.2912 & 0.2739 & 0.1934 \\
    \bottomrule
  \end{tabular}
\end{table}


We varied one hyperparameter at a time while keeping the remaining settings in
\cref{tab:hparams} fixed.
Each entry in \cref{tab:ablation} averages CLIP scores over the Burger, Cat, and Villa prompts.

\subsubsection{SDS: \texttt{guidance\_scale}}
\label{sec:abl_sds_gs}

\Cref{tab:ablation}(a) shows that the score rises from 0.2426 at guidance 7.5 to 0.2950 at 15,
with the highest observed score of 0.2991 at 25.
Larger scales reduce the score to 0.2879 at 50 and 0.2815 at 100.
In \cref{fig:abl_sds_gs}, guidance 7.5 produces blurred subjects, and the villa is barely visible
around the pool.
At 15, the subjects become more distinct, while 50 and 100 introduce increasingly dense background
textures, especially around the burger and cat.

Increasing CFG amplifies the difference between positive- and negative-condition noise predictions,
giving the noisy SDS updates a stronger direction toward the text prompt.
Excessively strong guidance can also exaggerate textures, consistent with the cluttered backgrounds
at 50 and 100.
Guidance 25 gives the best observed score among the tested values, compared with 7.5 commonly used
for ordinary sampling.
These results suggest that SDS benefits from stronger text guidance, but increasing it beyond a
moderate value does not necessarily improve alignment or visual quality.

\subsubsection{SDI: \texttt{inversion\_n\_steps}}
\label{sec:abl_sdi_steps}

One inversion step produces a low score of 0.1809 and largely unrecognizable images in
\cref{fig:abl_sdi_steps}, consistent with an overly coarse inversion.
Three steps recover recognizable subjects and raise the score to 0.2942.
The scores then vary non-monotonically: 0.2874 at 5 steps, 0.3006 at 10, and 0.2956 at 25.
More steps therefore do not guarantee better text alignment.
Smaller inversion steps provide a finer approximation to the predicted trajectory, but final image
quality also depends on target reconstruction and latent optimization.

Each additional inversion step requires another U-Net evaluation whenever the target is refreshed,
so runtime grows approximately linearly with the step count.
The trade-off is clear: 1 step is fast but fails to form recognizable subjects, whereas 25 steps
costs about 6.7 times as much as 3 steps for a score increase of only about 0.0013.
Three steps is therefore a lower-cost option with a similar CLIP score, while 10 steps gives the
highest observed score among the tested settings.

\subsubsection{SDI: \texttt{sdi\_update\_interval}}
\label{sec:abl_sdi_interval}

The highest observed score is 0.3006 at interval 25, while interval 1000 drops to 0.1934
(\cref{tab:ablation}(c)).

At the smallest tested interval of 25, the target is refreshed 40 times during 1000 iterations,
allowing it to follow the changing latent and incorporate progressively lower noise levels.
The subjects remain recognizable in \cref{fig:abl_sdi_interval}, and this setting gives the highest
observed score.
At the largest interval of 1000, the target is constructed only once, at the initial timestep 980.
All subsequent iterations fit that fixed target without receiving a new estimate at a lower noise
level, consistent with the severe artifacts and low score.
Larger intervals therefore hold potentially outdated targets for more iterations and skip more
scheduled timesteps.
The results favor repeated target updates over fitting a single initial target, although the
intermediate scores show that the effect is not strictly monotonic.

%-------------------------------------------------------------------------
\subsection{Hyperparameters}
\label{sec:hparams}

\input{Table/3_hparams}

\Cref{tab:hparams} lists the settings used for the main comparison.
The ablations vary only the parameter under study; all other settings remain fixed.
